\clearpage
\pagestyle{plain}
\chapter*{Digitális melléklet használati útmutató}
\addcontentsline{toc}{chapter}{Digitális melléklet használati útmutató}

A szakdolgozathoz tartozó digitális melléklet célja, hogy a dolgozatban bemutatott WaccApp rendszer forrásállományai, dokumentációs elemei, teszteléshez kapcsolódó anyagai és a dolgozat elektronikus változatai áttekinthető formában rendelkezésre álljanak. A melléklet nem a dolgozat főszövegének helyettesítésére szolgál, hanem annak gyakorlati kiegészítése. Segítségével ellenőrizhető, hogy a dolgozatban ismertetett frontendnézetek, API-kapcsolatok, kérdőíves folyamatok, médiakezelési megoldások és szűrési funkciók milyen projektállományokhoz kapcsolódnak.

A digitális melléklet neve: \path{WaccApp_digitalis_melleklet}. A melléklet könyvtárszerkezete a következő:

\begin{CodeBlock}
WaccApp_digitalis_melleklet/
|-- 01_dolgozat/
|   |-- szakdolgozat_Gal_Oliver_Istvan_ISXBF3.pdf
|   `-- latex_forras/
|       |-- dolgozat.tex
|       |-- chapters/
|       |-- cover/
|       |-- images/
|       `-- styles/
|
|-- 02_forraskod/
|   `-- waccapp/
|       |-- app.js
|       |-- questionnaire.js
|       |-- package.json
|       |-- package-lock.json
|       |-- .env.example
|       |-- kerdoivek/
|       |-- README.txt
|       `-- public/
|           |-- index.html
|           |-- conv.html
|           |-- filter.html
|           |-- messages.html
|           |-- sent-messages.html
|           |-- questionnaire.html
|           |-- template.html
|           |-- template-menu.html
|           |-- template-single.html
|           |-- login.html
|           |-- menu.html
|           |-- admin.html
|           |-- account.html
|           |-- privacy.html
|           |-- send.html
|           |-- questionnaires.json
|           |-- uploads/
|           `-- sent_media/
|
|-- 03_adatbazis/
|   `-- whatsapp_messages_teszt.db
|
|-- 04_dokumentacio/
|   |-- futtatasi_utmutato.txt
|   `-- felhasznaloi_utmutato.txt
|
`-- 05_teszteles/
    |-- tesztforgatokonyvek.pdf
    `-- tesztadatok_leirasa.txt
\end{CodeBlock}

A \path{01_dolgozat} mappa tartalmazza a szakdolgozat végleges PDF-változatát, valamint a dolgozat LaTeX-forrásállományait. A LaTeX-források között található a fő \path{dolgozat.tex} fájl, a fejezeteket tartalmazó \path{chapters} mappa, az előzékeket tartalmazó \path{cover} mappa, a képeket tartalmazó \path{images} mappa és a formázáshoz szükséges \path{styles} mappa. Ez a rész azért került a mellékletbe, hogy a dolgozat szerkesztett formában is visszakereshető legyen, ne csak a végleges PDF álljon rendelkezésre.

A \path{02_forraskod} mappa tartalmazza a WaccApp alkalmazás forrásállományait. Ebben található az \path{app.js} szerveroldali kapcsolódási pont, a \path{questionnaire.js} kérdőíves logikát tartalmazó állomány, valamint a függőségeket leíró \path{package.json} és \path{package-lock.json}. A mappában szereplő \path{.env.example} fájl a szükséges környezeti változók mintáját mutatja be. A tényleges konfigurációs értékek, hozzáférési adatok, API-tokenek és jelszavak biztonsági okokból nem kerültek bele a digitális mellékletbe, ezért ezek helyett csak a változók szerkezete látható.

A \path{02_forraskod/waccapp/public} mappa tartalmazza a felhasználó által közvetlenül elérhető frontendnézeteket. Ide tartozik többek között az \path{index.html}, amely az üzenetküldés, sablonküldés, kérdőívindítás, médiaküldés és időzítés központi felülete. A \path{conv.html} a kontaktalapú beszélgetésnézetet biztosítja, a \path{filter.html} a szűrési és egyszerű válaszösszesítési feladatokat támogatja, a \path{messages.html} és a \path{sent-messages.html} pedig a beérkező és elküldött üzenetek listaalapú áttekintéséhez kapcsolódik. A kérdőíves és sablonos szerkesztési feladatokat a \path{questionnaire.html}, a \path{template.html}, a \path{template-menu.html} és a \path{template-single.html} oldalak támogatják.

A \path{public/uploads} és a \path{public/sent_media} mappák a média- és fájlkezeléshez kapcsolódnak. Ezek a mappák a rendszer működésében a feltöltött, illetve elküldött médiaállományok kezeléséhez szükségesek. A digitális mellékletben ezek csak bemutatási vagy tesztelési célú állományokat tartalmaznak, mert a dolgozat célja nem éles kommunikációs tartalmak átadása, hanem a fejlesztett rendszer szerkezetének és működésének bemutatása. Emiatt éles kommunikációból származó fájlok, személyes adatok, hozzáférési adatok vagy tokeneket tartalmazó állományok nem kerültek a mellékletbe.

A \path{03_adatbazis} mappa a teszteléshez és bemutatáshoz használt SQLite-adatbázist tartalmazza. A mellékletben szereplő adatbázis tesztadatbázisként értelmezendő. Célja az, hogy a dolgozatban bemutatott üzenetküldési, visszakeresési, beszélgetésmegjelenítési és kérdőíves folyamatok ellenőrizhetők legyenek. Az adatbázis ezért nem éles felhasználói adatokat vagy valós ügyfélkommunikációt tartalmaz, hanem a rendszer működésének bemutatásához használható adatokat.

A \path{04_dokumentacio} mappa tartalmazza a futtatási és felhasználói útmutatót. A \path{futtatasi_utmutato.txt} röviden bemutatja a projekt indításának alaplépéseit, például a szükséges függőségek telepítését, a környezeti változók beállítását és a szerver indítását. A \path{felhasznaloi_utmutato.txt} a WaccApp főbb frontendnézeteinek használatát foglalja össze, külön kitérve az üzenetküldésre, a beszélgetésnézetre, a szűrésre, a kérdőívszerkesztésre és a sablonos folyamatokra.

Az \path{05_teszteles} mappa a manuális teszteléshez kapcsolódó anyagokat tartalmazza. A \path{tesztforgatokonyvek.pdf} a dolgozat mellékletében is bemutatott manuális tesztforgatókönyvek külön dokumentált változata lehet, míg a \path{tesztadatok_leirasa.txt} röviden ismerteti, hogy milyen tesztadatokkal történt a rendszer kipróbálása. Ezek az anyagok a dolgozat tesztelési fejezetében és mellékletében bemutatott ellenőrzési szempontokat egészítik ki.

A digitális melléklet összeállításakor fontos szempont volt, hogy a csomag kizárólag a dolgozat ellenőrzéséhez és a fejlesztett rendszer bemutatásához szükséges állományokat tartalmazza. Emiatt a mellékletben nem szerepelnek fejlesztői környezethez kötődő, automatikusan újragenerálható vagy biztonsági szempontból érzékeny fájlok. A \path{node_modules} mappa például nem része a mellékletnek, mert a szükséges függőségek a \path{package.json} és a \path{package-lock.json} alapján újratelepíthetők. A \path{.git} és a \path{.idea} mappák szintén nem kerültek bele, mivel ezek a verziókezeléshez, illetve a helyi fejlesztői környezethez kapcsolódnak, és nem szükségesek a rendszer szakdolgozati bemutatásához.

A valódi \path{.env} fájl helyett a mellékletben \path{.env.example} található. Ennek oka, hogy a tényleges \path{.env} fájl olyan konfigurációs adatokat tartalmazhat, amelyek hozzáférési tokenekhez, API-kulcsokhoz, session beállításokhoz vagy más érzékeny adatokhoz kapcsolódnak. A melléklet célja nem ezek átadása, hanem annak bemutatása, hogy a rendszer milyen típusú környezeti változókat használ. Ezért a mintafájl a szükséges változók nevét és szerkezetét mutatja be, de valódi titkos értékeket nem tartalmaz.

A projekt futtatásához a \path{02_forraskod/waccapp} mappában található forrásállományok használhatók. A függőségek a \path{package.json} alapján telepíthetők, a környezeti változók pedig a \path{.env.example} fájl alapján állíthatók be. A szükséges konfigurációs adatok megadása után a szerveroldali alkalmazás indítható, majd a frontendnézetek böngészőből érhetők el. A dolgozat frontendfókusza miatt a digitális melléklet legfontosabb szerepe az, hogy a bemutatott HTML-nézetek, kliensoldali működések, képernyőképek és API-kapcsolatok ellenőrizhető projektkörnyezetben is visszakereshetők legyenek.

A digitális melléklet szerkezete tehát azt a célt szolgálja, hogy a szakdolgozat elméleti és gyakorlati része egymáshoz kapcsolható legyen. A dolgozat főszövege bemutatja a WaccApp frontendjének tervezési és működési logikáját, a melléklet pedig biztosítja azokat az állományokat, amelyek alapján a megvalósított rendszer felépítése, forráskódja, tesztelése és futtatási környezete áttekinthető.